\documentclass[11pt,a4paper]{article}

\usepackage[a4paper,margin=1.55cm]{geometry}
\usepackage{fontspec}
\usepackage{xeCJK}
\usepackage{xcolor}
\usepackage{enumitem}
\usepackage{titlesec}
\usepackage{tabularx}
\usepackage[hidelinks]{hyperref}

\IfFontExistsTF{PingFang SC}{
  \setCJKmainfont{PingFang SC}
}{
  \IfFontExistsTF{Noto Sans CJK SC}{
    \setCJKmainfont{Noto Sans CJK SC}
  }{
    \IfFontExistsTF{Source Han Sans SC}{
      \setCJKmainfont{Source Han Sans SC}
    }{}
  }
}

\pagestyle{empty}
\setlength{\parindent}{0pt}
\setlength{\parskip}{0pt}
\linespread{1.08}
\XeTeXlinebreaklocale "zh"
\XeTeXlinebreakskip = 0pt plus 1pt

\setlist[itemize]{leftmargin=1.2em,itemsep=0.18em,topsep=0.18em,parsep=0em,partopsep=0em}

\titleformat{\section}{\large\bfseries}{}{0pt}{}[\titlerule]
\titlespacing*{\section}{0pt}{0.55em}{0.35em}

\titleformat{\subsection}{\bfseries}{}{0pt}{}
\titlespacing*{\subsection}{0pt}{0.35em}{0.15em}

\newcommand{\Header}[4]{
  {\LARGE\bfseries #1}\par
  {\normalsize #2}\par
  \vspace{0.25em}
  {\small #3}\par
  {\small #4}\par
  \vspace{0.6em}
}

\newcommand{\RoleLine}[3]{
  \begin{tabularx}{\textwidth}{@{}X r@{}}
    \textbf{#1} & #3 \\
  \end{tabularx}
  {\small #2}\par
}

\begin{document}

\Header
  {余胜民}
  {全栈软件交付工程师}
  {17600512027 · \href{mailto:susuyan@163.com}{susuyan@163.com} · 微信: susuyan\_dream · \href{https://github.com/susuyan}{github.com/susuyan}}
  {专注于端到端系统交付，使用 AI 驱动的开发工作流。过去一年让 Claude Code 和 Cursor 写了 40\% 的代码，然后告诉客户这本来就是计划的一部分。最擅长把模糊的需求变成可交付的产品。}

\section*{PROJECT}
\subsection*{PICO VR 采集平台}
\RoleLine{全栈交付顾问}{}{2025.11 -- Present}
\begin{itemize}
  \item 设计并交付跨平台 VR 数据采集系统（Unity/C\#/Node.js/React/Flutter）
  \item 实现高并发数据处理服务，吞吐量达 10,000+ 条传感器数据/秒
  \item 集成 PICO VR SDK 实现空间定位追踪，数据采集精度提升 35\%
  \item 提供端到端解决方案，项目周期缩短 30\%
\end{itemize}

\section*{EXPERIENCE}
\subsection*{iOS Engineer — 深圳海伦司企业管理有限公司}
\RoleLine{}{}{2021.04 -- 2025.09}
\begin{itemize}
  \item 架构并交付覆盖 500+ 门店的 ERP 和 OA 系统，运营效率提升 40\%
  \item 设计智能音乐播放系统，门店管理成本降低 60\%
  \item 使用 Docker 和 GitHub Actions 建立 CI/CD 流程，部署频率从每周提升到每日
  \item 建立代码评审和自动化测试体系，缺陷率降低 45\%
\end{itemize}
\subsection*{iOS Engineer - 北京鱼乐贝贝教育科技有限公司}
\RoleLine{}{}{2016.12 -- 2021.02}
\begin{itemize}
  \item 负责 iOS 客户端开发及小团队（3–4 人）技术推进，参与业务需求拆解、任务分配与版本迭代，支撑公司在全国近两千家门店的业务系统稳定运行。
  \item 主导将家长端与教师端老项目从 Objective-C 重构为 Swift，完成核心模块重构与架构优化，提升代码可维护性与开发效率，支撑后续功能快速迭代。
  \item 参与多端产品建设，独立开发子公司及幼儿园客户端应用，并主导 iPad 端「加盟宝」项目，落地门店地图、投资计算器及数据看板等核心功能，提升加盟商决策与运营效率。
\end{itemize}
\subsection*{Software Engineer — 北京亿维鸿科技有限公司}
\RoleLine{}{}{2015.07 -- 2016.11}
\begin{itemize}
  \item 负责多款 iOS 客户端应用的开发与维护，涵盖工具类产品（手机归属地工具、手机管家）及游戏类应用，参与功能开发、问题修复与版本迭代，逐步熟悉客户端开发流程及 App Store 上架规范。
\end{itemize}

\section*{技能}
{\small TypeScript · React · Python · Node.js · Docker · Claude Code · Cursor · Swift · PostgreSQL · Kubernetes · AWS · Flutter · Unity 3D · C\# · CI/CD · 系统设计 · 快速交付 · 全栈思维 · AI 工具链集成}

\section*{教育背景}
\textbf{洛阳理工学院} \;|\; 计算机应用技术

\end{document}
